\scanpage{121}
% archive foliation: 120
Si $S$ est de car. 0, alors $X$ est \inserted{(lisse/$S$ ??)} à fibres géom. connexes non vides, \unsure{on peut énoncer que} pour $s,t:X\longrightarrow X_0$ (fibré en droites projectives sur $S$) étales, tels que $G(t,s)=0$ i.e. $t^*(c_{X_0})=s^*(c_{X_0})$, il existe
\[
g\in\operatorname{Aut}(X_0)
\]
tel que $s=g\circ t$.
\margin{\uline{à vérifier !}}
Donc, pour $t$ fixé, $s$ variable, le morphisme
\[
f^{-1}(H_S)\backslash\mathcal P_X\longrightarrow\deleted{\unsure{\Omega}}\underline\omega_{X/S}^{\otimes2}\qquad s\longmapsto G(t,s)
\]
est \uline{injectif}, ou encore (intrinsèquement) le morphisme
\[
H\backslash\mathcal P_X\xrightarrow{c^X}\operatorname{Conn}_{X/S}
\]
défini par passage au quotient par $c^X$, est \uline{injectif}. \deleted{$\unsure{H\backslash}$} Mais bien sûr ces deux morphismes ne sont plus \deleted{bijectifs} surjectifs. Une conn. \unsure{affine|projective} qui est \deleted{\unsure{contenue}} dans l'image de $H\backslash\mathcal P_X$ est dite ``intégrable par fonctions algébriques''.

Si $S$ est de car. pos., alors l'injectivité n'est pas non plus valable non plus. Introduisons $X^{(p/S)}$ et
\diagram{\begin{tikzcd}
X\arrow[rr,"F_X"]\arrow[dr]&&X^{p/S}\arrow[dl]\\
&S&
\end{tikzcd}} (Frobenius)
On trouve que le \deleted{\unsure{groupe $H_{X/S}$}} \inserted{$H(X^{p/S})$}, opère \unsure{rég.|s.ég.} \uline{de façon évidente}, sur $\deleted{\unsure{\mathcal P_X}}\ \Gamma(X,\mathcal P_X)=\operatorname{Ét}_S(X,\mathbb P^1_S)$, laisse invariant l'expression $G(t,s)$, et le même \unsure{argumentation opère}, i.e. d'où encore
\[
F_X^{-1}(H_{X^{p/S}})\backslash\mathcal P_X\lhook\joinrel\longrightarrow\operatorname{Conn}_{X/S}.
\]
\margin{\unsure{si $2$ et $3$ sont inversibles, en car. $p>3$ !} Cf plus bas !}

Introduisons (\unsure{si $s,t$ sont} \deleted{deux} \inserted{\unsure{deux}} uniformisants) l'opérateur
\[
D_t^{(n)}f=\frac1{n!}\frac{d^nf}{dt^n}
\]
(qui a un sens sans restriction de car. nulle, dès que $X/S$ est diff. lisse de \deleted{\unsure{dim}} \inserted{dim rel.} ($=1$)), on trouve alors que (pour tout uniformisant)
\[
G(t,s)=12\left(-\frac{D_t^{(3)}s}{D_t^{(1)}s}+\left(\frac{D_t^{(2)}s}{D_t^{(1)}s}\right)^2\right)=12G_0(t,s)
\]
où
\[
G_0(t,s)=-\frac{D_t^{(3)}s}{\unsure{D_t^{(1)}}s}+\left(\frac{D_t^{(2)}s}{\unsure{D_t^{(1)}}s}\right)^2.
\]
Cela montre que si $12\cdot1_S\deleted{\unsure{\ne}}=0$, alors toutes les connexions projectives de la forme $c_t$ sont \uline{égales}, donc il y a une conn. proj. canonique sur $X$ rel. à $S$, à laquelle on peut référer toutes les autres. Cela montre aussi (sans restriction de car. telle $12\cdot1_S=0$ !) que le torseur $\deleted{\unsure{\mathcal T}}\operatorname{Conn}_{X/S}$ est déduit, moyennant le chgt de groupe d'op.
\[
\underline\omega_{X/S}^{\otimes2}\xrightarrow{12\cdot\mathrm{id}}\underline\omega_{X/S}^{\otimes2},
\]
d'un torseur plus fin $\operatorname{Conn}_{0\,X/S}$ \inserted{sous $\underline\omega_{X/S}^{\otimes2}$}, avec un morphisme
\[
\mathcal P_{X/S}\xrightarrow{c_0^X}\operatorname{Conn}_{0\,X/S}
\]
tel que
\[
c_0^X(s)-c_0^X(t)=G_0(t,s).
\]
Bien entendu, les assertions plus haut sur l'expression $G(t,s)$ et les connexions
